\bookchapter{Глава 65. Проект 1 — моя книга учёта (локальная версия)}

\chapterintro{Добро пожаловать в ваш первый «серьёзный проект». Эта книга учёта не выходит в сеть, но устроена по полной программе: записывает транзакции, показывает историю, сводит статистику по месяцам, умеет удалять и править. Сделав её, вы целиком пройдёте весь цикл «данные логика интерфейс». Не спешите срисовывать код — сперва подумайте, потом делайте: так мы договаривались с самого начала.}

\booksection{Мыслим требованиями: сперва задайте пять вопросов}

\begin{center}
\begin{tabularx}{\linewidth}{|X|X|}
\hline
\textbf{Элемент} & \textbf{Ответ} \tabularnewline
\hline
Who (для кого) & Для себя, можно и для родных и друзей \tabularnewline
What (что делаем) & Записывать доходы и расходы, смотреть историю, сводить по месяцам, считать статистику по категориям \tabularnewline
Why (зачем) & Тренировка + понять, куда утекают деньги \tabularnewline
How (как делаем) & Swift + SwiftUI + SwiftData + MVVM, полностью локально, без бэкенда \tabularnewline
Constraint (ограничения) & Уложиться в неделю; минимум iOS 17 (порог SwiftData) \tabularnewline
\hline
\end{tabularx}
\end{center}

\begin{center}\small\color{gray}Пять элементов требований книги учёта\end{center}

Затем разбейте работу на мелкие шаги, где «после каждого шага уже можно пользоваться»: (1) слой данных на SwiftData (можно сохранить запись); (2) страница-список (можно смотреть историю); (3) окно ввода (можно добавить запись); (4) удаление и правка; (5) карточка месячной сводки. После каждого шага приложение остаётся «работающей вещью».

\booksection{Слой данных: даём записям дом}

Сначала опишите, как выглядит одна запись, — поля нужны достаточные, но без жадности:

\begin{lstlisting}[language=Swift]
import SwiftData

@Model
class Record {
    var type: String      // "income" доход | "expense" расход
    var amount: Double    // сумма (в юанях)
    var category: String  // категория: еда, транспорт, зарплата…
    var note: String      // заметка (необязательно)
    var date: Date        // дата записи

    init(type: String, amount: Double, category: String, note: String = "", date: Date = Date()) {
        self.type = type
        self.amount = amount
        self.category = category
        self.note = note
        self.date = date
    }
}
\end{lstlisting}

\begin{tcbinfo}{Совет: почему сумма — это Double?}
Строгие финансовые приложения берут Decimal (точная десятичная арифметика), чтобы избежать ошибок плавающей точки (вечная проблема 0.1+0.2$\neq$0.3). Для учебного проекта хватает Double; если однажды будете делать настоящий финансовый продукт, не забудьте перейти на Decimal или «Int в копейках» — на эти грабли наступили бессчётные предшественники.

\end{tcbinfo}

\booksection{MVVM: интерфейс и логика живут раздельно}

Дать интерфейсу право напрямую трогать базу? Работать будет, но код интерфейса разрастётся в длинную неопрятную простыню. Архитектура MVVM разводит всех по обязанностям: View (интерфейс) только рисует, ViewModel (логический центр) распоряжается данными, Model (SwiftData) ведает хранением.

View (интерфейс SwiftUI)

ViewModel (логический центр на @Observable)

Model (слой данных SwiftData)

\begin{center}\small\color{gray}Рисунок 4: трёхслойная структура MVVM — данные текут в одну сторону, слои не заходят на чужую территорию\end{center}

\begin{lstlisting}[language=Swift]
import Observation

@Observable
class RecordViewModel {
    var records: [Record] = []
    private let context: ModelContext

    init(context: ModelContext) { self.context = context }

    func load() {
        // фильтрация + сортировка через #Predicate в SwiftData
        let descriptor = FetchDescriptor<Record>(
            sortBy: [SortDescriptor(\.date, order: .reverse)]
        )
        records = (try? context.fetch(descriptor)) ?? []
    }

    func add(type: String, amount: Double, category: String, note: String) {
        let r = Record(type: type, amount: amount, category: category, note: note)
        context.insert(r)
        try? context.save()
        load()                       // перезагружаем — интерфейс обновится
    }

    func delete(_ record: Record) {
        context.delete(record)
        try? context.save()
        load()
    }
}
\end{lstlisting}

В интерфейсе создаём @State private var vm = RecordViewModel(...) и рисуем список прямо из vm.records — данные изменились, интерфейс обновился сам.

\booksection{Интерфейс: страница-список + окно ввода}

\begin{lstlisting}[language=Swift]
struct RecordListView: View {
    @State private var vm: RecordViewModel
    @State private var showAdd = false

    var body: some View {
        NavigationStack {
            List {
                Section("Баланс за месяц") {          // сводная карточка
                    SummaryRow(records: vm.records)
                }
                Section("Операции") {
                    ForEach(vm.records) { r in
                        RecordRow(record: r)
                    }
                    .onDelete { idx in
                        idx.forEach { vm.delete(vm.records[$0]) }
                    }
                }
            }
            .navigationTitle("Мои записи расходов")
            .toolbar {
                Button { showAdd = true } label: {
                    Image(systemName: "plus")
                }
            }
            .sheet(isPresented: $showAdd) {
                AddRecordView(vm: vm)          // диалог ввода
            }
        }
    }
}
\end{lstlisting}

Главное в окне ввода: у поля суммы клавиатура .keyboardType(.decimalPad) (цифровая), доход/расход — сегментный переключатель (Picker), категория — выбор в меню, по кнопке «Сохранить» вызываем vm.add(...) и закрываем окно через dismiss().

\begin{tcbwarn}{Внимание: проверку ввода не упрощаем}
Поле суммы обязательно проверять: пустое значение, отрицательные числа, недопустимый ввод вроде «abc» — всё перехватываем (проверяем через Double(text), при неудаче он вернёт nil). Одна из самых частых причин падений приложений — «посчитать пользовательский ввод послушным ребёнком». Запомните: любой ввод пользователя досматривайте как потенциального врага.

\end{tcbwarn}

\booksection{Правка, удаление и сводка}

\begin{itemize}
  \item Удаление: .onDelete у List — смахнули влево и удалили + окно подтверждения (случайное удаление — то, что пользователи учётных приложений ненавидят сильнее всего);
  \item Правка: нажатие на строку открывает окно редактирования с уже подставленными старыми значениями (поправить одно слово проще, чем вводить всё заново);
  \item Сводка: через filter + reduce считаем «доход / расход / остаток за месяц» и показываем карточкой вверху:
\end{itemize}

\begin{lstlisting}[language=Swift]
func monthSummary(_ records: [Record]) -> (income: Double, expense: Double) {
    let cal = Calendar.current
    let monthRecords = records.filter {
        cal.isDate($0.date, equalTo: Date(), toGranularity: .month)
    }
    let income = monthRecords
        .filter { $0.type == "income" }
        .reduce(0) { $0 + $1.amount }
    let expense = monthRecords
        .filter { $0.type == "expense" }
        .reduce(0) { $0 + $1.amount }
    return (income, expense)
}
\end{lstlisting}

\booksection{Приёмка: три уровня проверки + два задания со звёздочкой}

Проект готов — принимайте по трём уровням «функции границы архитектура»: записывается ли, удаляется ли, остаются ли данные после перезапуска приложения (функции); не упадёт ли при вводе суммы 0 или отрицательного числа (границы); не перемешаны ли интерфейс с данными, нет ли жёстко вбитых значений (архитектура).

Когда базовая версия прошла приёмку, добавьте себе два «блюда повыше классом»:

\begin{enumerate}
  \item Статистика по категориям: добавьте вкладку «Статистика» и рисуйте траты по категориям в виде круговой или столбчатой диаграммы через Swift Charts — официальная библиотека графиков Apple: пара строк кода, и диаграмма готова;
  \item Синхронизация iCloud: чтобы включить iCloud у SwiftData, нужно лишь несколько настроек — книга учёта сама синхронизируется между iPhone и iPad, и проект мгновенно становится «по-взрослому».
\end{enumerate}

\begin{tcbpractice}{Практика: закончите его за эту неделю}
Действуйте в порядке «слой данных ViewModel страница-список окно ввода карточка сводки», по шагу в день. Застряли — отдайте ошибку в первозданном виде поисковику или ИИ (помните часть 7 «Практика программирования с ИИ»? Описывая проблему, ясно формулируйте «ожидаемое поведение, фактическое поведение, шаги воспроизведения»). Когда закончите, ведите в нём учёт неделю — пользоваться приложением, сделанным своими руками, приятно, и это чувство успеха подтолкнёт вас идти дальше.

\end{tcbpractice}

\begin{tcbsummary}{Краткие итоги главы}
\begin{itemize}
  \item Пять элементов требований разбить на мелкие шаги, где «после каждого шага всё работает», и лишь потом браться за код
  \item Слой данных: модель @Model + контейнер modelContainer + запрос FetchDescriptor
  \item Разделение в MVVM: View только рисует, @Observable-ViewModel ведает логикой, SwiftData — хранением
  \item Интерфейс: List + Section + toolbar + окно ввода через sheet
  \item Трёхуровневая приёмка (функции / границы / архитектура) + задания со звёздочкой (Swift Charts, синхронизация iCloud)
\end{itemize}

\end{tcbsummary}
